\documentclass[11pt]{article}

% Latin Modern + T1: scalable outlines (microtype expansion needs them) and
% accented author names in the bibliography (e.g. B\l{}\k{a}dek).
\usepackage{lmodern}
\usepackage[T1]{fontenc}
\usepackage[margin=0.9in]{geometry}
\usepackage{amsmath, amssymb, mathtools}
\usepackage{amsthm}
\usepackage{algorithm, algpseudocode}
\usepackage{microtype}
\usepackage{enumitem}
\usepackage{booktabs}
\usepackage{longtable}
\usepackage{array}
\usepackage{xcolor}
\definecolor{linkblue}{RGB}{0, 70, 140}
\usepackage[numbers,sort&compress]{natbib}
\usepackage[colorlinks=true, urlcolor=linkblue, linkcolor=linkblue, citecolor=linkblue]{hyperref}
\usepackage[nameinlink,capitalize,noabbrev]{cleveref}

\setlength{\parskip}{3pt}
\setlength{\parindent}{0pt}
\setlength{\emergencystretch}{2em}
\pagestyle{plain}

% math shorthands, matching methodology.tex
\newcommand{\calD}{\mathcal{D}}
\newcommand{\calH}{\mathcal{H}}
\newcommand{\calS}{\mathcal{S}}
\newcommand{\calI}{\mathcal{I}}
\newcommand{\calP}{\mathcal{P}}
\newcommand{\bbH}{\mathbb{H}}
\newcommand{\bbR}{\mathbb{R}}
\newcommand{\Lhat}{\hat{L}}
\DeclareMathOperator*{\argmin}{arg\,min}

% Hyperlinked pointers into the student skeleton (GitHub blob/tree URLs).
\newcommand{\repoBlob}{https://github.com/krichelj/ByteBoost-Shape-Constrained-Boosted-SR/blob/main/}
\newcommand{\repoTree}{https://github.com/krichelj/ByteBoost-Shape-Constrained-Boosted-SR/tree/main/}
\newcommand{\pyFile}[1]{\href{\repoBlob#1}{\nolinkurl{#1}}}
\newcommand{\pyDir}[1]{\href{\repoTree#1}{\nolinkurl{#1/}}}
\newcommand{\codeRef}[1]{%
  \par\nobreak\noindent
  {\footnotesize\textit{Code skeleton:}~#1}\par\smallskip}

\theoremstyle{plain}
\newtheorem{proposition}{Proposition}
\newtheorem{theorem}{Theorem}
\newtheorem{lemma}{Lemma}
\theoremstyle{remark}
\newtheorem*{remark*}{Remark}

\begin{document}

% ======================================================================
%  ByteBoost 2026 - Project Description
% ======================================================================

\begin{center}
 {\Large\bfseries ByteBoost Workshop: Project Description}\\[6pt]
    {\large\bfseries Shape-Constrained Boosted Symbolic Regression\\
    for Interpretable Neural Scaling Laws: A Cross-Testbed Investigation}
\end{center}

\hrule
\vspace{6pt}

\section{Field of science}\label{sec:field}
Computer Science: Machine Learning and Artificial Intelligence, specifically the foundations of large language models (neural scaling laws) and interpretable machine learning via symbolic regression. Secondary field: High-Performance Computing and computer architecture (cross-platform performance characterization).

\section{Abstract / summary}\label{sec:abstract}

Training a modern AI language model can cost millions of dollars, so researchers use scaling laws to predict performance from model size and training data before committing resources. These equations guide major hardware and budget decisions, but their mathematical form is usually chosen manually and then fitted to measured results.

Our project uses symbolic regression to discover the scaling-law equation directly from data. Symbolic regression searches many candidate formulas and returns an interpretable equation rather than a black-box model. Unconstrained searches, however, often produce formulas that fit existing measurements but extrapolate unrealistically, predicting negative error or worse performance with more data.

We therefore require every candidate formula to satisfy admissibility axioms: larger models and additional data cannot increase error; improvements must diminish with scale; error must approach a positive irreducible floor; excess loss must decay as a power law; and the formula must remain smoothly well-defined down to the smallest scales of interest. The discovery method is soft shape-constrained search: a genetic-programming fitness \emph{penalizes} continuum axiom violations via interval-arithmetic scores, drives them toward zero, and bounds any remaining shortfall (\cref{sec:soft,prop:soft,sec:algorithm,sec:axioms,sec:stage-admiss}).

The formula is built incrementally from a standard scaling law, adding one correction at a time while preserving all guarantees (\cref{prop:chinchilla-admiss,prop:stage-iff,thm:guarantee}). We also prove that the corrections improve fit without changing the irreducible error floor or the asymptotic rate of improvement whenever the soft scores vanish (\cref{thm:guarantee,prop:soft}).

The project has two computational bottlenecks. Model training will use Cerebras CS-3 systems on PSC's Neocortex platform; new validation losses will be compared to already-published loss curves from an existing scaling grid on Hugging Face (\cref{sec:testbeds,sec:baselines}). Formula discovery is a large CPU-bound search dominated by interval-arithmetic certification. Because it is highly parallel and relies mainly on scalar or imperfectly vectorized Python (with optional JAX acceleration), Stony Brook's high-core-count AMA27 AmpereOne cluster is well suited to the workload (\cref{sec:testbeds,sec:certificates}).

Deliverables will include a certified scaling-law formula fitted to our training data, an open-source discovery library, and a concise methods report (\cref{sec:deliverables}). An initial set of trained models is already public (\cref{sec:datasets}).

\section{Key requirements}\label{sec:requirements}

\subsection{Testbed hardware (requested)}\label{sec:testbeds}
\codeRef{\pyDir{src/modeling/training}, \pyDir{src/systems/hpc}, \pyDir{src/search}}
\begin{itemize}[nosep, leftmargin=1.4em]
 \item \textbf{Neocortex (PSC):} Cerebras CS-3 wafer-scale systems for the model-pretraining grid in \cref{sec:models,sec:datasets}
 (\url{https://www.psc.edu/resources/neocortex/}).
 Each CS-3 is built around a Wafer Scale Engine~3 (WSE-3) with on-chip SRAM and extreme memory bandwidth, so a model in our~$N$ range can stay on a single wafer without multi-GPU model parallelism.
 We use Neocortex only for pretraining; the resulting validation losses are compared to the already-published loss curves of \cref{sec:baselines} (not to re-running those models on student-accessible GPUs).
 \item \textbf{AMA27 (Stony Brook Supercomputing Center / ACCESS):} Arm-based AmpereOne A192-32M cluster for the CPU-bound symbolic-regression search of \cref{sec:boosting} and the certification inner loop of \cref{sec:certificates}.
 Per the ACCESS resource description, phase~1 deploys 312 single-socket compute nodes, each with a 192-core AmpereOne A192-32M CPU and 384\,GB DDR5, linked by 200\,Gbps NDR InfiniBand, with a 5\,PB GPFS filesystem (\url{https://support.access-ci.org/affinity-groups/ama27}).
 AMA27 is the Stony Brook Arm testbed featured in ByteBoost~2.0 (in place of the retired Ookami A64FX system).
 We do \emph{not} rely on A64FX-style 512-bit SVE or HBM: the GP/IA workload is population-parallel and largely scalar Python DualInterval evaluation, so the useful AMA27 properties are Arm ISA diversity, high core counts per node, and dense node count for concurrent candidate evaluation.
\end{itemize}

\subsection{Existing scaling-law baselines (public artifacts)}\label{sec:baselines}
\codeRef{\pyFile{src/scaling/data/scaling_dataset.py}, \pyFile{src/modeling/training/trainer.py}}
An already-public \texttt{colinear\_\allowbreak scaling\_\allowbreak models} collection on Hugging Face~\citep{kricheli2025colinear} provides trained LLM checkpoints and the associated validation-loss curves from a prior scaling grid.
Those artifacts are the workshop baseline: students compare \emph{losses}~$L(\mathbf{h})$ (and the scaling laws fit to them), not the original training hardware, which is neither part of this project nor available to participants.
New Neocortex pretraining runs should emit losses on a comparable protocol so they can be plotted and analyzed against that public loss inventory (\cref{sec:datasets,sec:models}).

\subsection{Software and tools}\label{sec:software}
\codeRef{\pyDir{src/modeling/training}, \pyDir{src/search}, \pyDir{src/constraints/certificates}}
PyTorch~2.x for pretraining (distributed data-parallel training; optional mixed precision), with \texttt{uv} for environment management and Weights \& Biases for experiment logging~\citep{biewald2020wandb}.
Symbolic search uses \texttt{gplearn}~\citep{stephens2015gplearn} with a custom fitness that adds interval-arithmetic axiom penalties---the sole discovery method of \cref{sec:soft}; an optional PySR backend provides an unconstrained alternate engine for ablation.
Certificates are evaluated with a pure-Python forward-mode \emph{DualInterval} interval-arithmetic layer (\cref{sec:certificates}); JAX with custom JVPs is an optional acceleration path, not required.

\subsection{Models}\label{sec:models}
\codeRef{\pyFile{src/modeling/models/config.py}, \pyFile{src/modeling/models/transformer.py}}
Decoder-only transformer language models (a compact post-LN causal stack suitable for scaling grids), spanning roughly 10M--1.4B parameters, across several tokens-per-parameter ratios ($M = D/N$).
Workshop ports may adopt LLaMA-style blocks~\citep{touvron2023llama}, FlashAttention, and $\mu$P~/~$\mu$Transfer~\citep{yang2022mup}; the core method only requires consistent $(N,D,L)$ measurements in the Step Law~\citep{li2025steplaw} / overtraining~\citep{gadre2024overtraining} regimes.

\subsection{Datasets}\label{sec:datasets}
\codeRef{\pyFile{src/scaling/data/scaling_dataset.py}, \pyFile{src/scaling/data/corpora.py}}
Pretraining corpora used to generate new grid points include Wikipedia~\citep{wikimedia2023wikipedia}, RedPajama~\citep{together2023redpajama}, and related open text mixtures (optional FineWeb-Edu~\citep{penedo2024finewebedu} for workshop extensions).
An already-public \texttt{colinear\_\allowbreak scaling\_\allowbreak models} collection on Hugging Face~\citep{kricheli2025colinear} supplies trained checkpoints and loss curves for the symbolic-regression track and for loss-level comparison with new Neocortex runs (\cref{sec:baselines}); the original training hardware is out of scope.

\section{Method summary (technical)}\label{sec:method}
\codeRef{\pyDir{src/scaling}, \pyDir{src/constraints}, \pyDir{src/search}; map: \cref{sec:code-map}}
Symbols used below are collected in \cref{sec:notation}; the student skeleton layout is summarized in \cref{sec:code-map}.

\subsection{Setup}\label{sec:setup}
\codeRef{\pyFile{src/scaling/setup/configuration.py}, \pyFile{src/scaling/setup/constants.py}}
Let $N$ be the parameter count, $D$ the training-token count, and
$\calH = \{h_1, \ldots, h_m\}$ a finite set of remaining hyperparameters
(in the concrete 4D grids used here, typically width-multiplier / weight-decay style axes and learning rate); see also \cref{sec:notation}.
For each coordinate $h \in \{N, D\} \cup \calH$, let $H_h$ be its finite set of
admissible values and $\phi_h \colon H_h \to \bbR_{>0}$ a strictly monotone
preprocessing map; write $\phi_i \coloneqq \phi_{h_i}$ for brevity.
In the search implementation, genetic-programming features are the logarithms
$\log\phi_h(h)$ (so stage corrections are fit in log-coordinates and mapped back
to raw-scale derivatives by the chain rule; \cref{sec:certificates}).
The configuration space is the product
$\bbH \coloneqq \prod_{h \in \{N,D\}\cup\calH} H_h$,
with typical element $\mathbf{h} \in \bbH$ (so $N, D, h_i$ denote the corresponding
coordinates of~$\mathbf{h}$).
Let $L \colon \bbH \to \bbR$ be validation loss; we seek $\Lhat \approx L$.
We are given a labeled dataset
$\calD = \{(\mathbf{h}_\ell, L(\mathbf{h}_\ell))\}_{\ell=1}^{n} \subset \bbH \times \bbR$.
For each scale variable $x \in \{N, D\}$ set
$x_{\min} \coloneqq \min H_x > 0$ and
$x_{\max} \coloneqq \max H_x$,
treat $x$ as continuous on
$[x_{\min}, \infty)$, and take all derivatives with the other coordinates held fixed.
Write $\widetilde{\bbH}$ for the corresponding continuum domain: each scale ranges over
$[x_{\min},\infty)$ while non-scale coordinates still range over their finite~$H_h$,
so $\bbH \subset \widetilde{\bbH}$.
The labeled loss~$L$ is observed only on~$\bbH$; the approximant~$\Lhat$ is evaluated
on~$\widetilde{\bbH}$ when stating shape constraints (\cref{sec:axioms}).

\subsection{Admissibility axioms}\label{sec:axioms}
\codeRef{\pyFile{src/constraints/axioms/admissibility.py}}
$\Lhat$ is \emph{admissible} ($\Lhat \in \calS$) if it satisfies~(A1)--(A5) below.
Where an axiom mentions a scale variable $x$, it is imposed for each
$x \in \{N, D\}$ with the other coordinates held fixed.
Although the labeled grid~$\bbH$ is finite, $\Lhat$ is constrained as a real function on
the continuum domain~$\widetilde{\bbH}$ of \cref{sec:setup}, so that the derivatives
in~(A1)--(A2) and the limits in~(A3)--(A4) are well-defined.

\begin{enumerate}[nosep, leftmargin=2.2em, label=\textup{(A\arabic*)}]
    \item \textbf{Monotone decrease}~\citep{hestness2017scaling,kaplan2020scaling}.
    $\dfrac{\partial \Lhat}{\partial x} \leq 0$ for all $x \geq x_{\min}$.
    \item \textbf{Diminishing returns}~\citep{hestness2017scaling,kaplan2020scaling}.
    $\dfrac{\partial^2 \Lhat}{\partial x^2} \geq 0$ for all $x \geq x_{\min}$.
    (Together with~(A1), this is the usual diminishing-returns shape: further
    increases in~$x$ yield weakly smaller loss reductions.)
    \item \textbf{Irreducible loss}~\citep{hestness2017scaling,hoffmann2022chinchilla}.
    There is a positive \emph{joint floor} $L_\infty$ and, along each axis, a
    \emph{marginal floor} $L^{\infty}_{x}$:
    \begin{enumerate}[nosep, leftmargin=1.4em, label=\textup{(\roman*)}]
        \item \emph{Joint floor.}
        There exists $L_\infty > 0$ such that
        $\Lhat(\mathbf{h}) > L_\infty$ for every $\mathbf{h} \in \widetilde{\bbH}$.
        \item \emph{Marginal floor.}
        For each $x \in \{N,D\}$ and every fixed choice of the other
        coordinates, the single-axis limit
        \begin{equation}\label{eq:margfloor-bb}
            L^{\infty}_{x}
            \;\coloneqq\;
            \lim_{t \to \infty}
            \Lhat(\mathbf{h})\big|_{x = t}
        \end{equation}
        exists and is finite, and satisfies
        $L^{\infty}_{x} \geq L_\infty$.
        (The value $L^{\infty}_{x}$ depends on the frozen coordinates; the notation
        suppresses that dependence.)
    \end{enumerate}
    In the stage-0 Chinchilla law, $L_\infty$ equals the joint limit as all scales
    $\to\infty$ (namely~$E$), while a single-axis marginal is in general strictly
    larger.
    Joint-floor positivity~$L_\infty>0$ already forces~$\Lhat>0$ on~$\widetilde{\bbH}$, so a
    separate positivity axiom is unnecessary.
    \item \textbf{Asymptotic power-law decay}~\citep{hestness2017scaling,kaplan2020scaling,hoffmann2022chinchilla}.
    Relative to the marginal floor~\eqref{eq:margfloor-bb},
    \begin{equation}\label{eq:powerlaw-bb}
        \lim_{t\to\infty}
        \frac{\log\bigl(\Lhat(\mathbf{h})\big|_{x=t}
            - L^{\infty}_{x}\bigr)}{\log t}
        \;=\; c_x
        \qquad\text{for some } c_x < 0.
    \end{equation}
    The excess $\Lhat(\mathbf{h})|_{x=t}-L^{\infty}_{x}$ is eventually positive (so the
    logarithm is defined) whenever the limit in~(A3)(ii) is approached from above,
    which follows from~(A1) when the floor is not attained at finite~$t$.
    \item \textbf{Graceful saturation}~\citep{kronberger2022shape}.
    $\Lhat$ is twice continuously differentiable ($C^2$) in~$x$ on $[x_{\min}, \infty)$
    ---so that the second derivative in~(A2) exists and is continuous---and
    $\Lhat(x_{\min}, \cdot)$ is finite.
    (Expression-tree corrections used later are in fact $C^\infty$ on $(0,\infty)$.
    The inequality $\Lhat(x_{\min},\cdot)>L_\infty$ is already required by~(A3)(i).)
\end{enumerate}

These encode physical priors on loss surfaces that unconstrained symbolic regression routinely violates.
In particular,~(A4) must use $L^{\infty}_{x}$ rather than $L_\infty$: along a single axis the excess over the joint floor need not vanish, so a power-law rate would not be identifiable against $L_\infty$.

\paragraph{Where the axioms come from.}
(A1)--(A4) are not new modeling choices; each is the formal statement of an empirical regularity already reported in the neural scaling-law literature, and we adopt them from that literature rather than positing them.
Power-law scaling of generalization error in model and data size is the central finding of \citet{hestness2017scaling} and \citet{kaplan2020scaling}, and is what~(A4) asserts; the monotone decrease~(A1) and the diminishing returns~(A2) are the shape that a decaying power law has.
\citet{hestness2017scaling} further decompose measured learning curves into a small-data region, a power-law region, and an \emph{irreducible-error region} bounded below by Bayes error---which is the floor~$L_\infty$ of~(A3).
\citet{hoffmann2022chinchilla} make that floor explicit in the parametric form $\Lhat(N,D) = E + A N^{-\alpha} + B D^{-\beta}$, in which~$E$ is read as the loss of an ideal generative process on the data distribution, i.e.\ the entropy of natural text.
Our stage-0 law~\eqref{eq:chinchilla-bb} is exactly that form generalized to additional hyperparameter axes, so the constants $L_\infty = E$ and $c_N^{(0)} = -\alpha$, $c_D^{(0)} = -\beta$ of \cref{prop:chinchilla-admiss} inherit their published meaning directly---which is what makes preserving them across boosting stages (\cref{thm:guarantee}) worth proving.
Related joint $(N,D)$ functional forms fitted for extrapolation appear in \citet{rosenfeld2020constructive} and \citet{alabdulmohsin2022revisiting}.
By contrast,~(A5) is a regularity condition rather than an empirical claim: it is the smoothness that derivative-based shape constraints need in order to be stated at all~\citep{kronberger2022shape}.

\paragraph{Scope of~(A1)--(A2).}
Monotonicity and convexity are a modeling restriction, not a universal law.
\citet{caballero2023broken} document \emph{non-monotonic} scaling behavior---double descent, sharp inflections---that no monotone form can express.
We impose~(A1)--(A2) because this project targets the smooth, monotonically decreasing regime of LLM pretraining validation loss on the grids of \cref{sec:datasets}.
A testbed that exhibited breaks or double descent would require the axiom set to be relaxed; detecting such a failure is itself a legitimate outcome here, since the soft scores of \cref{sec:soft} \emph{measure} the admissibility gap rather than assume it away.

\subsection{Boosted construction}\label{sec:boosting}
\codeRef{\pyDir{src/search/baselines}, \pyDir{src/search/residuals}, \pyDir{src/search/expression}}
The ensemble is built stagewise,
$\Lhat^{(j)} = \Lhat^{(j-1)} + \Lhat_j$, targeting an \emph{admissibility invariant}:
$\Lhat^{(j)} \in \calS$ for every $j = 0, \ldots, K$ (\cref{sec:axioms}),
enforced by the shape-constrained search of \cref{sec:soft} whenever $V=0$ with finite DualInterval enclosures (\cref{prop:soft,thm:guarantee}).
Throughout, $\Lhat_j$ (subscript) is the stage-$j$ correction and
$\Lhat^{(j)}$ (superscript) is the cumulative ensemble through stage~$j$,
so in particular $\Lhat^{(0)} = \Lhat_0$.
The stagewise-additive construction is gradient boosting in the sense of \citet{friedman2001greedy}, with the shape-constrained search of \cref{sec:soft} in place of an unconstrained base learner.
Stage~$0$ is a generalized Chinchilla law~\citep{hoffmann2022chinchilla}, fit by nonlinear least squares with
positive parameter constraints.
Concrete instances include a 2D law in $(N,D)$ and a 4D law that adds two
hyperparameter axes (e.g.\ weight decay and learning rate). The 4D hyperparameter
terms may be power laws of the same form as in~\eqref{eq:chinchilla-bb}, or---matching
U-shaped loss in $\mathrm{lr}$/$\mathrm{wd}$---quadratic penalties in
$\log\phi_{\mathrm{lr}}(\mathrm{lr})$ and $\log\phi_{\mathrm{wd}}(\mathrm{wd})$
(shape axioms still act only on~$N$ and~$D$):
\begin{equation}\label{eq:chinchilla-bb}
    \Lhat^{(0)} \;=\; \Lhat_0 \;=\;
    \frac{A}{N^{\alpha}} + \frac{B}{D^{\beta}}
    + \sum_{i=1}^{m} \frac{C_i}{\phi_i(h_i)^{\gamma_i}} + E,
    \qquad A, B, C_i, E, \alpha, \beta, \gamma_i > 0 .
\end{equation}
\begin{proposition}[Stage-0 admissibility]\label{prop:chinchilla-admiss}
The law~\eqref{eq:chinchilla-bb} with all coefficients and exponents strictly positive
lies in~$\calS$, with joint floor $L_\infty = E$ and asymptotic exponents
$c_N^{(0)} = -\alpha$, $c_D^{(0)} = -\beta$ (the rates $c_x$ in~(A4) of \cref{sec:axioms} along $N$ and~$D$).
\end{proposition}
(See \cref{sec:chinchilla-proof} for the proof.)
Each later stage $j \geq 1$ fits Huber pseudo-residuals~\citep{huber1964robust,friedman2001greedy} of the current ensemble,
with clipping threshold
\begin{equation}\label{eq:delta-bb}
    \delta_j
    \;\coloneqq\;
    \mathrm{median}_{\ell=1,\ldots,n}
    \bigl|L(\mathbf{h}_\ell) - \Lhat^{(j-1)}(\mathbf{h}_\ell)\bigr|
\end{equation}
so that outlier runs cannot dominate:
\begin{equation}\label{eq:pseudoresid-bb}
    \tilde{r}_j^{(\ell)}
    \;\coloneqq\;
    \begin{cases}
        L(\mathbf{h}_\ell) - \Lhat^{(j-1)}(\mathbf{h}_\ell)
            & \text{if }
              \bigl|L(\mathbf{h}_\ell) - \Lhat^{(j-1)}(\mathbf{h}_\ell)\bigr|
              \leq \delta_j, \\[3pt]
        \delta_j \cdot
          \mathrm{sign}\bigl(L(\mathbf{h}_\ell)
                        - \Lhat^{(j-1)}(\mathbf{h}_\ell)\bigr)
            & \text{otherwise}.
    \end{cases}
\end{equation}
Corrections are sought by genetic-programming search over expression trees with
primary operators $\{+, -, \times, \div\} \cup \{\mathrm{pow}_p : p \in \calP\}$
for a finite set $\calP \subset \bbR \setminus \{0\}$
(where $\mathrm{pow}_p(z) = |z|^p$ for $z > 0$) and bounded depth;
optional unary helpers such as $\mathrm{inv}$/$\mathrm{log}$/$\mathrm{sqrt}$ may be enabled for ablations.
Trees are evaluated on \emph{log-features} $\log\phi_h(h)$; raw-scale derivatives
needed by~(A1)--(A2) are recovered by the chain rule (\cref{sec:certificates}).

\subsection{Stagewise admissibility}\label{sec:stage-admiss}\label{sec:certificates}
\codeRef{\pyFile{src/constraints/axioms/stage_conditions.py}, \pyFile{src/constraints/certificates/interval.py}, \pyFile{src/constraints/certificates/ia_eval.py}, \pyFile{src/constraints/certificates/ia_certificates.py}}
Checking that the full ensemble $\Lhat^{(j)}$ obeys~(A1)--(A5) of \cref{sec:axioms} is a continuum statement on~$\widetilde{\bbH}$ and along each scale half-line.
The next claim reduces that check to conditions on the increment alone.
\begin{equation}\label{eq:stage-admiss-bb}
\begin{gathered}
\text{(i) } \frac{\partial \Lhat_j}{\partial x}
   \leq -\frac{\partial \Lhat^{(j-1)}}{\partial x}
   \quad\text{for all } x \geq x_{\min},
\\[2pt]
\text{(ii) } \frac{\partial^2 \Lhat_j}{\partial x^2}
   \geq -\frac{\partial^2 \Lhat^{(j-1)}}{\partial x^2}
   \quad\text{for all } x \geq x_{\min},
\\[2pt]
\text{(iii) } \Lhat_j(\mathbf{h}) > -\Lhat^{(j-1)}(\mathbf{h})
   \quad\text{for all } \mathbf{h} \in \widetilde{\bbH},
\\[2pt]
\text{(iv) } \Lhat^{(j-1)}(\mathbf{h}) + \Lhat_j(\mathbf{h}) > L_\infty
   \quad\text{for all } \mathbf{h} \in \widetilde{\bbH},
\\[2pt]
\text{(v) } \Lhat_j(\mathbf{h})\big|_{x=t} = O\bigl(t^{c_{x,j}}\bigr)
   \text{ as } t \to \infty
   \text{ for some } c_{x,j} < c_x^{(0)},
\\[2pt]
\text{(vi) } \Lhat_j \in C^\infty
   \text{ on } [x_{\min}, x_{\max}]
   \text{ with }
   \bigl|\Lhat_j(x_{\min}, \cdot)\bigr| < \infty.
\end{gathered}
\end{equation}
When $L_\infty > 0$,~(iv) already implies~(iii); both are retained so positivity of the
ensemble and preservation of the joint floor remain explicit.
\begin{proposition}[Stagewise equivalence]\label{prop:stage-iff}
Given $\Lhat^{(j-1)} \in \calS$ with joint floor $L_\infty$, the update
$\Lhat^{(j)} = \Lhat^{(j-1)} + \Lhat_j$ lies in~$\calS$ with the same joint floor
if and only if $\Lhat_j$ satisfies~\eqref{eq:stage-admiss-bb}(i)--(vi) for each $x \in \{N,D\}$.
\end{proposition}
(See \cref{sec:stage-iff-proof} for the proof.)
Conditions~(i)--(iv) and~(vi) are still quantified over continua ((i)--(ii) on the half-line $x\geq x_{\min}$; (iii)--(iv) on~$\widetilde{\bbH}$), and~(v) is an asymptotic statement: none of these can be decided by sampling finitely many points of~$\bbH$.
They become machine-checkable because every GP proposal~$g$ is a concrete finite expression tree built from~$\{+,-,\times,\div\}\cup\{\mathrm{pow}_p:p\in\calP\}$. Write $F \coloneqq \Lhat^{(j-1)} + g$. Two complementary certificate constructions supply the continuous checks used by the soft fitness of \cref{sec:soft}: interval enclosures establish~(i)--(iv) and~(vi) on the compact box~$\calI_x$ covering the observed grid, while the structural check decides~(v) exactly (and thereby controls the $x\to\infty$ tail that the box does not cover).

\smallskip\noindent
\emph{(a) Sound interval-arithmetic checks} (conditions~(i)--(iv),~(vi) on a compact slice of~$\widetilde{\bbH}$).
Over the compact box $\calI_x \coloneqq [x_{\min}, x_{\max}]$ (other coordinates ranging over their finite $H_h$), the product of boxes is a compact continuum subset of~$\widetilde{\bbH}$ covering the observed scale range; the half-line $x>x_{\max}$ is handled by the structural test~(b) below.
Forward-mode automatic differentiation composed with interval arithmetic~\citep{moore2009interval,kronberger2022shape} evaluates~$g$ bottom-up on interval operands and returns guaranteed enclosures
$\underline{f} \leq f \leq \overline{f}$ on~$\calI_x$ for $f \in \{F,\, \partial F/\partial x,\, \partial^2 F/\partial x^2\}$, in time linear in the tree size.
(The second-derivative enclosure is what makes~(A2) / stage~(ii) checkable;~(A5) asks only for~$C^2$, while the operator set of \cref{sec:boosting} in fact yields~$C^\infty$ on $(0,\infty)$.)
When~$g$ is represented on log-features, DualInterval evaluation is performed in log-coordinates and first/second derivatives are mapped to the raw axis~$x$ by the chain rule before applying~\eqref{eq:interval-bb}.
Those enclosures are exactly the inputs to the soft violation scores of \cref{sec:soft}: writing $v_{\mathrm{mono}}=\max(0,\overline{\partial F/\partial x})$ and likewise for the other gaps, the sufficient inequalities
\begin{equation}\label{eq:interval-bb}
    \overline{\partial F / \partial x} \leq 0,
    \quad
    \underline{\partial^2 F / \partial x^2} \geq 0,
    \quad
    \underline{F} > L_\infty ,
    \quad
    \overline{F} < \infty
\end{equation}
hold for each $x \in \{N,D\}$ precisely when the corresponding soft gaps vanish (and $\overline{F}<\infty$).
Because the enclosures are sound, vanishing gaps are a \emph{sufficient} certificate that the continuum conditions hold \emph{everywhere} on~$\calI_x$, not merely at sampled points of~$\bbH$ (\cref{lem:ia-sound,sec:cert-lemmas}):
$\overline{\partial F/\partial x} \leq 0$ forces~$\partial F/\partial x \leq 0$ on~$\calI_x$ (hence~(i) on the box), and likewise for~(ii);
$\underline{F} > L_\infty > 0$ forces~(iv) on the continuum configurations covered by the product of boxes, and hence also~(iii) (ensemble positivity);
$\overline{F} < \infty$ forces the finiteness half of~(vi), while~$C^2$ (in fact~$C^\infty$) smoothness of~$g$ on~$\calI_x$ follows from the operator set and discharges~(A5) on the box.
The implication is one-sided: inflated enclosures can produce positive soft scores for a truly admissible update, but every~$g$ with vanishing gaps on~\eqref{eq:interval-bb} is guaranteed to preserve admissibility on~$\calI_x$.
Evaluating DualInterval enclosures (and the soft scores built from them) on each GP candidate is the inner-loop hot spot, the direct target of the AMA27 performance work.

\smallskip\noindent
\emph{(b) Exact structural checks} (condition~(v), and prerequisites for it).
Define the asymptotic exponent $\mathrm{ord}(g,x)$ recursively on the tree: constants and leaves other than~$x$ contribute~$0$; $\mathrm{pow}_p(x)$ contributes~$p$; products add exponents; quotients subtract; sums take the maximum.
Then~(v) holds with $c_{x,j} = \mathrm{ord}(g,x)$ precisely when $\mathrm{ord}(g,x) < c_x^{(0)}$ for each $x \in \{N,D\}$ (\cref{lem:ord,sec:cert-lemmas}).
We additionally require that each scale variable $x \in \{N,D\}$ appears as a leaf of~$g$ (otherwise the exponent in~$x$ is undefined or vacuously~$0$).
Both tests are exact, finite walks over a finite tree, with no floating-point search over~$\widetilde{\bbH}$.

\subsection{Shape-constrained search}\label{sec:soft}
\codeRef{\pyFile{src/search/soft/violations.py}, \pyFile{src/search/soft/fitness.py}, \pyFile{src/search/soft/backends.py}}
The discovery method---and the only search used by Algorithm~\ref{alg:boosting-bb}---is soft interval-arithmetic search: per-axiom violation scores on the candidate ensemble $F \coloneqq \Lhat^{(j-1)}+g$ are folded into the \texttt{gplearn} fitness (\cref{sec:software}), steering the genetic program toward admissible corrections.
(The student package is \pyDir{src/search/soft}; there is no separate reject-filter path.)
Using the interval enclosures of \cref{sec:stage-admiss}(a) and the structural quantities of~(b), define the soft scores of \cref{sec:notation}:
\begin{equation}\label{eq:violations-bb}
\begin{aligned}
    v_{\mathrm{mono}}(g)
      &= \max_{x\in\{N,D\}}
         \max\Bigl(0,\; \overline{\partial F/\partial x}\Bigr),
    &\quad&\text{(monotonicity gap)},
    \\
    v_{\mathrm{conv}}(g)
      &= \max_{x\in\{N,D\}}
         \max\Bigl(0,\; -\underline{\partial^2 F/\partial x^2}\Bigr),
    &\quad&\text{(convexity gap)},
    \\
    v_{\mathrm{irred}}(g)
      &= \max_{x\in\{N,D\}}
         \max\bigl(0,\; L_\infty - \underline{F}\bigr),
    &\quad&\text{(floor gap)},
    \\
    v_{\mathrm{decay}}(g)
      &= \max_{x\in\{N,D\}}
         \max\bigl(0,\;
           \mathrm{ord}(g,x) - c_x^{(0)}
           + \varepsilon_{\mathrm{decay}}\bigr),
    &\quad&\text{(power-law gap)},
    \\
    v_{\mathrm{leaf}}(g)
      &= \bigl|\{x\in\{N,D\}
           : x \notin \mathrm{leaves}(g)\}\bigr|,
    &\quad&\text{(missing scale leaves)}.
\end{aligned}
\end{equation}
Here $\varepsilon_{\mathrm{decay}} > 0$ is a fixed margin so that $v_{\mathrm{decay}}(g)=0$ forces the strict inequality $\mathrm{ord}(g,x) \leq c_x^{(0)}-\varepsilon_{\mathrm{decay}} < c_x^{(0)}$ required by~\eqref{eq:stage-admiss-bb}(v); $\mathrm{leaves}(g)$ is the set of variable leaves of the expression tree~$g$.
Every score is nonnegative, and $v_{\mathrm{mono}}=v_{\mathrm{conv}}=v_{\mathrm{irred}}=v_{\mathrm{leaf}}=0$ iff the corresponding interval or structural test of~\eqref{eq:interval-bb} /~(b) succeeds (for $v_{\mathrm{decay}}$, zero is sufficient but slightly stronger than~(v) because of the margin).
Writing $a$ for an axiom index ranging over $\{\mathrm{mono},\mathrm{conv},\mathrm{irred},\mathrm{decay},\mathrm{leaf}\}$ and $\lambda_a > 0$ for its penalty weight, the aggregate violation is $V(g)=\sum_a v_a(g)$.
Stage~$j$ then minimizes the penalized fitness
\begin{equation}\label{eq:penalized-bb}
    F_j(g) \;=\;
    \underbrace{\frac{1}{n} \sum_{\ell=1}^{n}
      \bigl(\tilde{r}_j^{(\ell)} - g(\mathbf{h}_\ell)\bigr)^2}_{\text{data fit}}
    \;+\;
    \underbrace{\sum_a \lambda_a \, v_a(g)^2}_{\text{axiom penalty}} ,
\end{equation}
where $F_j$ is a scalar fitness (distinct from the ensemble map~$F$ above), the first term is mean squared error of~$g$ against the Huber pseudo-residuals~\eqref{eq:pseudoresid-bb}, and the second term penalizes axiom gaps (squared so the penalty gradient vanishes at zero violation).
Candidates with vanishing soft scores (and finite DualInterval enclosures) therefore incur no penalty pressure; note that $v_{\mathrm{decay}}=0$ is slightly stronger than~\eqref{eq:stage-admiss-bb}(v) because of the margin~$\varepsilon_{\mathrm{decay}}$.

\paragraph{Where the penalties come from.}
Constraining a symbolic-regression model through bounds on its image and its partial derivatives, and computing those bounds with interval arithmetic, is the shape-constrained symbolic regression of \citet{kronberger2022shape}; \citet{bladek2019formal} attack the same problem with SMT solvers and counterexample-driven GP instead.
Our $v_{\mathrm{mono}}$, $v_{\mathrm{conv}}$, and $v_{\mathrm{irred}}$ are that literature's monotonicity, convexity, and image-bound constraints, specialized to~(A1)--(A3) of \cref{sec:axioms} and evaluated on the ensemble $F = \Lhat^{(j-1)} + g$ rather than on the candidate~$g$ alone---the modification that makes them meaningful stagewise (\cref{prop:stage-iff}).
Where we depart from \citet{kronberger2022shape} is the \emph{handling} of a violation: they enforce constraints as a hard feasibility filter, discarding infeasible candidates during selection or segregating them into a second population, whereas~\eqref{eq:penalized-bb} adds $\sum_a \lambda_a v_a(g)^2$ to the fitness and lets violations shrink continuously.
That soft, penalty-based treatment follows \citet{martinek2024penalized}, who minimize shape-constraint violations inside the objective rather than applying them only as a post-hoc selection filter, and who report the clearest gains in the data-scarce regime---which is the regime a scaling grid of expensive pretraining runs puts us in.
The one-sidedness recorded after~\eqref{eq:interval-bb}---inflated enclosures may penalize a candidate that is in fact admissible, but never certify one that is not---is the \emph{pessimistic} constraint evaluation compared against optimistic sampling-based evaluation by \citet{haider2022optimistic}; we take the pessimistic side deliberately, because \cref{thm:guarantee} needs soundness rather than tightness.
The remaining two scores have no counterpart in that literature: $v_{\mathrm{decay}}$ and $v_{\mathrm{leaf}}$ encode the asymptotic-rate axiom~(A4), which is specific to scaling laws, and they do so through the structural exponent $\mathrm{ord}(g,x)$ of \cref{sec:stage-admiss}(b) on the half-line $x > x_{\max}$ rather than through an interval bound on a compact box.
\begin{proposition}[Zero-penalty implies certificates]\label{prop:soft}
If $V(g)=0$ at every stage \emph{and} DualInterval evaluation returns finite upper
enclosures~$\overline{F}<\infty$ on each~$\calI_x$ (the finiteness half of~\eqref{eq:interval-bb},
discharged in practice by the $C^\infty$ operator set of \cref{sec:boosting}), then certificates~(a)--(b) hold and the boosting guarantee of \cref{thm:guarantee} applies exactly.
The scores $v_{\mathrm{mono}}$, $v_{\mathrm{conv}}$, $v_{\mathrm{irred}}$ upper-bound the corresponding true pointwise violations on~$\calI_x$ (while $v_{\mathrm{decay}}$ and $v_{\mathrm{leaf}}$ are exact).
\end{proposition}
(See \cref{sec:soft-proof} for the proof.)
The soft scores therefore yield a \emph{measured} admissibility gap, in contrast to unconstrained search ($\lambda_a=0$), which can fit the grid while leaving axiom violations unpenalized.
Soft scores act as IA proxies for~(A1)--(A4) (with~(A5) discharged by finiteness of the DualInterval enclosure / $C^2$---in practice $C^\infty$---operator set).
Comparing soft-penalty search ($\lambda_a>0$) against that unconstrained baseline is one of the empirical questions the project settles.

\subsection{Algorithm}\label{sec:algorithm}
\codeRef{\pyFile{src/search/boosting/algorithm.py}, \pyDir{src/search/boosting}}
Putting the pieces together, Algorithm~\ref{alg:boosting-bb} builds the ensemble of \cref{sec:boosting}.
At every stage the correction is the minimizer of the shape-constrained fitness~\eqref{eq:penalized-bb} from \cref{sec:soft}.
\begin{algorithm}[H]
\caption{Shape-constrained boosted symbolic regression}
\label{alg:boosting-bb}
\begin{algorithmic}[1]
\Require Dataset $\calD = \{(\mathbf{h}_\ell, L(\mathbf{h}_\ell))\}_{\ell=1}^{n}$,
         stages $K$
\Ensure Ensemble $\Lhat^{(K)}$ (aims for $V=0$ with finite DualInterval enclosures at every stage)
\State Fit $\Lhat_0$ via~\eqref{eq:chinchilla-bb} on $\calD$;
       set $\Lhat^{(0)} \gets \Lhat_0$
\For{$j = 1, \ldots, K$}
    \State Compute $\delta_j$ via~\eqref{eq:delta-bb}
    \State Compute $\tilde{r}_j^{(\ell)}$ for all $\ell=1,\ldots,n$
           via~\eqref{eq:pseudoresid-bb}
    \State Find $\Lhat_j$ by minimizing~\eqref{eq:penalized-bb}
    \State $\Lhat^{(j)} \gets \Lhat^{(j-1)} + \Lhat_j$
\EndFor
\State \Return $\Lhat^{(K)}$
\end{algorithmic}
\end{algorithm}

\subsection{Guarantee}\label{sec:guarantee}
\codeRef{\pyFile{src/constraints/guarantee/invariants.py}}
The soft search of \cref{sec:soft} drives each correction toward vanishing violation scores.
When $V(\Lhat_j)=0$ with finite DualInterval enclosures at every stage (\cref{prop:soft}), certificates~(a)--(b) of \cref{sec:stage-admiss} hold, and the ensemble preserves admissibility, the joint floor, and the stage-0 rates:
\begin{theorem}[Boosting guarantee]\label{thm:guarantee}
Assume $\Lhat_0$ is the admissible stage-0 law of \cref{prop:chinchilla-admiss} (so $L_\infty = E$).
If every later correction $\Lhat_j$ ($j\geq 1$) has $V(\Lhat_j)=0$ with finite DualInterval enclosures---equivalently, satisfies certificates~(a)--(b) of \cref{sec:stage-admiss} (\cref{prop:soft})---then for all $j$: $\Lhat^{(j)} \in \calS$; writing $L_\infty^{(j)}$ for the joint floor of~$\Lhat^{(j)}$, one has $L_\infty^{(j)} = E$; and the asymptotic exponents are preserved,
\begin{equation}\label{eq:guarantee-bb}
    \lim_{t \to \infty}
    \frac{\log\bigl(\Lhat^{(j)}(\mathbf{h})\big|_{x=t}
        - L^{\infty}_{x}\bigr)}{\log t}
    \;=\; c_x^{(0)},
\end{equation}
where $L^{\infty}_{x}$ is the marginal floor of $\Lhat^{(j)}$ along~$x$ (equal to that of~$\Lhat_0$).
\end{theorem}
Boosting therefore improves fit \emph{without} corrupting the interpretable constants researchers actually quote (proofs in \cref{sec:guarantee-proof,sec:soft-proof}).

\section{Additional information}\label{sec:additional}

\subsection{Collaborators welcome}\label{sec:collaborators}
\codeRef{\pyDir{src/modeling} (modeling); \pyDir{src/search}, \pyDir{src/systems/hpc} (search)}
The project splits into two fairly independent tracks, so a teammate could own either one.
The \emph{modeling} track ports and benchmarks transformer pretraining on Neocortex's Cerebras CS-3 systems (\cref{sec:testbeds}) and compares the resulting validation losses to the public Hugging Face loss curves of \cref{sec:baselines} (useful background: PyTorch, or curiosity about the Cerebras software stack).
The \emph{search} track ports the genetic-programming symbolic-regression loop of \cref{sec:boosting,sec:algorithm,sec:soft} to AMA27 (\cref{sec:testbeds}) and profiles the DualInterval certification inner loop of \cref{sec:certificates} on AmpereOne (useful background: C/C++ or Python performance work, or an interest in Arm/HPC profiling); a self-contained slice of it is benchmarking soft-penalty search against $\lambda_a=0$ (\cref{sec:soft}).
No prior scaling-law expertise is needed for either; the main thing is an appetite for making code run fast on unusual hardware.
The student skeleton paths for each track are listed in \cref{sec:code-map}.

\subsection{Deliverables}\label{sec:deliverables}
\codeRef{\pyFile{src/systems/pipeline/run_discovery.py}}
(1) An admissible boosted scaling law (\cref{sec:axioms,sec:guarantee}) fit on the collected training grid (\cref{sec:datasets}); (2) a public release of the shape-constrained symbolic-regression library with the DualInterval soft-score / certification stack of \cref{sec:certificates,sec:soft}; and (3) a short paper on the method of \cref{sec:method}.
Notation, the description$\leftrightarrow$skeleton map, and the formal proofs appear in \cref{sec:notation,sec:code-map,sec:proofs}.

\clearpage
\renewcommand{\refname}{References}
\renewcommand{\bibsection}{\section{\refname}}
\bibliographystyle{unsrtnat}
\bibliography{byteboost_refs}

\clearpage
\appendix

\section{Notation}\label{sec:notation}
Symbols used in \cref{sec:method} (and pointed to from the student stubs of \cref{sec:code-map}).
Formal claims from \cref{sec:method} are proved in \cref{sec:proofs}.
Violation scores~$v_a$, aggregate~$V$, and fitness~$F_j$ belong to the sole discovery method of \cref{sec:soft}.
\begin{longtable}{@{}>{$}l<{$}p{0.72\textwidth}@{}}
\toprule
\textbf{Symbol} & \textbf{Meaning} \\
\midrule
\endfirsthead
\toprule
\textbf{Symbol} & \textbf{Meaning} \\
\midrule
\endhead
\bottomrule
\endfoot
N, D
  & Parameter count and training-token count (scale variables). \\
M
  & Tokens-per-parameter ratio $M = D/N$ (\cref{sec:models}). \\
\calH = \{h_1,\ldots,h_m\}
  & Remaining hyperparameters; $m = |\calH|$. \\
i
  & Hyperparameter index, $i=1,\ldots,m$ (distinct from
    data-point index~$\ell$). \\
H_h
  & Finite set of admissible values for coordinate $h \in \{N,D\}\cup\calH$. \\
\phi_h,\; \phi_i
  & Strictly monotone preprocessing $\phi_h \colon H_h \to \bbR_{>0}$;
    $\phi_i \coloneqq \phi_{h_i}$. \\
\bbH
  & Discrete configuration space $\prod_{h} H_h$ (labeled grid). \\
\widetilde{\bbH}
  & Continuum domain of \cref{sec:setup,sec:axioms}: each scale ranges over
    $[x_{\min},\infty)$, other coordinates over finite~$H_h$;
    $\bbH \subset \widetilde{\bbH}$. \\
\mathbf{h}
  & Typical element of $\bbH$ or $\widetilde{\bbH}$ (coordinates $N,D,h_i$). \\
L,\; \Lhat
  & True validation loss and its approximation. \\
\calD,\; n,\; \ell
  & Labeled dataset
    $\{(\mathbf{h}_\ell, L(\mathbf{h}_\ell))\}_{\ell=1}^{n}$;
    $n = |\calD|$; data-point index~$\ell$. \\
x
  & Generic scale variable, $x \in \{N,D\}$. \\
t
  & Dummy limit variable along a scale axis
    (e.g.\ $\Lhat(\mathbf{h})|_{x=t}$). \\
x_{\min},\; x_{\max}
  & $\min H_x$ and $\max H_x$. \\
\calS
  & Set of admissible scaling laws (\cref{sec:axioms}). \\
L_\infty,\; L_\infty^{(j)}
  & Joint irreducible-loss floor; joint floor of~$\Lhat^{(j)}$
    ($L_\infty^{(j)}=E$ when $V=0$ with finite enclosures). \\
L^{\infty}_{x}
  & Marginal floor along axis~$x$~\eqref{eq:margfloor-bb}
    (depends on the frozen coordinates). \\
c_x,\; c_x^{(0)}
  & Asymptotic power-law exponent along~$x$~(A4);
    stage-0 value $c_N^{(0)}=-\alpha$, $c_D^{(0)}=-\beta$. \\
j,\; K
  & Boosting stage index; number of correction stages after
    stage~$0$ (\cref{sec:algorithm}, Algorithm~\ref{alg:boosting-bb}). \\
\Lhat_j,\; \Lhat^{(j)}
  & Stage-$j$ correction (subscript) and cumulative ensemble through
    stage~$j$ (superscript); $\Lhat^{(0)}=\Lhat_0$. \\
A,B,C_i,E,\alpha,\beta,\gamma_i
  & Positive coefficients/exponents of the stage-0 Chinchilla
    law~\eqref{eq:chinchilla-bb}; joint floor $L_\infty=E$. \\
\delta_j
  & Huber clipping threshold at stage~$j$~\eqref{eq:delta-bb}. \\
\tilde{r}_j^{(\ell)}
  & Huber pseudo-residual for data point~$\ell$ at
    stage~$j$~\eqref{eq:pseudoresid-bb}. \\
\calP,\; p
  & Finite set of admissible powers
    $\calP\subset\bbR\setminus\{0\}$; generic element~$p$. \\
\mathrm{pow}_p
  & Power primitive $\mathrm{pow}_p(z)=|z|^p$ for $z>0$. \\
g,\; F
  & Candidate expression tree; ensemble-plus-candidate
    $F=\Lhat^{(j-1)}+g$. \\
c_{x,j},\; \mathrm{ord}(g,x)
  & Asymptotic exponent of correction/candidate~$g$ along~$x$;
    $c_{x,j}=\mathrm{ord}(g,x)$ in~\eqref{eq:stage-admiss-bb}(v). \\
\calI_x
  & Compact certification box $[x_{\min},x_{\max}]$ along scale~$x$
    (written $I_x$ in the skeleton);
    with other coordinates over finite~$H_h$, the product of boxes is a
    compact continuum slice of~$\widetilde{\bbH}$ covering the observed
    grid (\cref{sec:stage-admiss}); the tail $x>x_{\max}$ is controlled
    by~$\mathrm{ord}$. \\
\underline{f},\; \overline{f}
  & Interval-arithmetic lower/upper enclosures of~$f$ on~$\calI_x$
    (global over that box for fixed~$x$; soft scores take
    $\max_{x\in\{N,D\}}$ of the per-axis gaps).
    $\underline{F}>L_\infty$ certifies~(A3) on the box (hence ensemble positivity);
    $\overline{F}<\infty$ with the smooth operator set certifies~(A5). \\
a,\; v_a(g),\; V(g)
  & Soft-search axiom index $a\in\{\mathrm{mono},\mathrm{conv},\mathrm{irred},
    \mathrm{decay},\mathrm{leaf}\}$; per-axiom violation score;
    aggregate $V=\sum_a v_a$~\eqref{eq:violations-bb}. \\
\varepsilon_{\mathrm{decay}}
  & Positive margin in $v_{\mathrm{decay}}$ so zero score implies
    strict power-law decay~\eqref{eq:violations-bb}. \\
\mathrm{leaves}(g)
  & Set of variable leaves of expression tree~$g$. \\
\lambda_a
  & Positive penalty weight for axiom~$a$. \\
F_j(g)
  & Shape-constrained penalized fitness at stage~$j$~\eqref{eq:penalized-bb}
    (distinct from~$F$). \\
\end{longtable}

\section{Code skeleton map}\label{sec:code-map}
The repository ships a student implementation skeleton under \pyDir{src}, grouped into five packages (\pyFile{src/README.md}).
Blue links open the corresponding paths on GitHub (\texttt{main}).
Each technical subsection in the main text also points to those files via the \textit{Code skeleton} lines; symbols used there are listed in \cref{sec:notation}.
The stage-0 baseline used by the boosting stubs is \cref{prop:chinchilla-admiss}; proofs of all formal claims are in \cref{sec:proofs}.
The table follows the soft-only method order of \cref{sec:method}: DualInterval enclosures feed the sole discovery package \pyDir{src/search/soft} (\cref{sec:soft}), which Algorithm~\ref{alg:boosting-bb} calls at every stage; there is no hard reject-filter package.

{\small
\setlength{\tabcolsep}{4pt}
\begin{longtable}{@{}>{\raggedright\arraybackslash}p{0.28\textwidth}
  l
  >{\raggedright\arraybackslash}p{0.42\textwidth}@{}}
\toprule
\textbf{Description} & \textbf{Package} & \textbf{Primary paths} \\
\midrule
\endfirsthead
\toprule
\textbf{Description} & \textbf{Package} & \textbf{Primary paths} \\
\midrule
\endhead
\bottomrule
\endfoot
\cref{sec:setup} & \pyDir{src/scaling} & \pyFile{src/scaling/setup/configuration.py}, \pyFile{src/scaling/setup/constants.py} \\
\cref{sec:datasets} & \pyDir{src/scaling} & \pyFile{src/scaling/data/scaling_dataset.py}, \pyFile{src/scaling/data/corpora.py} \\
\cref{sec:axioms} & \pyDir{src/constraints} & \pyFile{src/constraints/axioms/admissibility.py} \\
\cref{sec:boosting} & \pyDir{src/search} & \pyDir{src/search/baselines}, \pyDir{src/search/residuals}, \pyDir{src/search/expression} \\
\cref{sec:stage-admiss,sec:certificates} & \pyDir{src/constraints} & \pyFile{src/constraints/axioms/stage_conditions.py}, \pyFile{src/constraints/certificates/interval.py}, \pyFile{src/constraints/certificates/ia_eval.py}, \pyFile{src/constraints/certificates/ia_certificates.py} \\
\cref{sec:soft} & \pyDir{src/search} & \pyFile{src/search/soft/violations.py}, \pyFile{src/search/soft/fitness.py}, \pyFile{src/search/soft/backends.py} \\
\cref{sec:algorithm}, Alg.~\ref{alg:boosting-bb} & \pyDir{src/search} & \pyFile{src/search/boosting/algorithm.py}, \pyDir{src/search/boosting} \\
\cref{sec:guarantee} & \pyDir{src/constraints} & \pyFile{src/constraints/guarantee/invariants.py} \\
\cref{sec:software} (optional JAX) & \pyDir{src/constraints} & \pyFile{src/constraints/certificates/jax_backend.py} \\
\cref{sec:baselines} & \pyDir{src/scaling} & \pyFile{src/scaling/data/scaling_dataset.py} \\
\cref{sec:models} & \pyDir{src/modeling} & \pyFile{src/modeling/models/config.py}, \pyFile{src/modeling/models/transformer.py} \\
\cref{sec:testbeds} (Neocortex) & \pyDir{src/modeling} & \pyDir{src/modeling/training} \\
\cref{sec:testbeds} (AMA27/CPU) & \pyDir{src/systems} & \pyDir{src/systems/hpc} \\
\cref{sec:deliverables} & \pyDir{src/systems} & \pyFile{src/systems/pipeline/run_discovery.py} \\
\end{longtable}
}


\section{Proofs}\label{sec:proofs}
Proofs follow the dependency order of \cref{sec:method}: stagewise equivalence, then the certificate lemmas, then zero-penalty implies certificates, then the boosting guarantee.
Stage-0 admissibility (\cref{prop:chinchilla-admiss}) is independent of the boosting loop and is proved last.

\subsection{Stagewise equivalence}\label{sec:stage-iff-proof}
\begin{proof}[Proof of \cref{prop:stage-iff}]
Write $G\coloneqq\Lhat^{(j-1)}$ and $g\coloneqq\Lhat_j$, so $\Lhat^{(j)}=G+g$.
Throughout, derivatives along $x\in\{N,D\}$ hold the other coordinates fixed.

\paragraph{(A1)--(A2).}
$\partial(G+g)/\partial x=\partial G/\partial x+\partial g/\partial x$ and likewise for second derivatives.
Since $G\in\calS$, one has $\partial G/\partial x\leq 0$ and $\partial^2 G/\partial x^2\geq 0$ on $[x_{\min},\infty)$.
Thus $\partial(G+g)/\partial x\leq 0$ on that half-line if and only if~(i) holds, and $\partial^2(G+g)/\partial x^2\geq 0$ if and only if~(ii) holds.

\paragraph{Joint floor~(A3) and ensemble positivity.}
The joint-floor inequality $G+g>L_\infty$ on $\widetilde{\bbH}$ is exactly~(iv).
(Here $L_\infty$ is the joint floor already certified for~$G$;~(iv) keeps it for the update.
When $L_\infty>0$,~(iv) also yields ensemble positivity~$G+g>0$, which is~(iii).)

\paragraph{(A5).}
If $G$ is $C^2$ on $[x_{\min},\infty)$ and $g$ is $C^\infty$ on $[x_{\min},x_{\max}]$ with finite boundary values~(vi), then $G+g$ is $C^2$ on the certification box.
On the half-line beyond $x_{\max}$, the same $C^2$ regularity follows once~$g$ is realized by a finite expression tree from $\{+,-,\times,\div\}\cup\{\mathrm{pow}_p\}$ (each primitive is $C^\infty$ on $(0,\infty)$), which is the setting of \cref{sec:boosting}.
Finiteness of $(G+g)(x_{\min},\cdot)$ is the restriction of~(vi) together with~$G\in\calS$; the bound $(G+g)(x_{\min},\cdot)>L_\infty$ is the restriction of~(iv).

\paragraph{(A3) marginal floors and~(A4).}
Because $G\in\calS$, the marginal floor $L^{\infty}_{x}(G)$ exists and the excess $G|_{x=t}-L^{\infty}_{x}(G)$ is asymptotic to a constant times $t^{c_x^{(0)}}$ with $c_x^{(0)}<0$.
Condition~(v) gives $g|_{x=t}=O(t^{c_{x,j}})$ with $c_{x,j}<c_x^{(0)}$, so $g|_{x=t}=o\bigl(t^{c_x^{(0)}}\bigr)$.
Therefore
\[
  \lim_{t\to\infty}(G+g)\big|_{x=t}
  \;=\;
  L^{\infty}_{x}(G),
\]
the marginal floor is unchanged, and
\[
  \lim_{t\to\infty}
  \frac{\log\bigl((G+g)|_{x=t}-L^{\infty}_{x}(G)\bigr)}{\log t}
  \;=\; c_x^{(0)}.
\]
Conversely, if $G+g\in\calS$ with the same joint floor and the same asymptotic rate~$c_x^{(0)}$, the difference $g=(G+g)-G$ must be $o(t^{c_x^{(0)}})$ along each axis (else the leading excess exponent of $G+g$ would differ from that of~$G$), which is~(v); the remaining items~(i)--(iv) and~(vi) are the rearrangements already used above.

Thus $G+g\in\calS$ with joint floor $L_\infty$ if and only if~\eqref{eq:stage-admiss-bb} holds.
\end{proof}

\subsection{Interval and structural certificates}\label{sec:cert-lemmas}
\begin{lemma}[Soundness of the interval tests]\label{lem:ia-sound}
Suppose forward-mode DualInterval evaluation returns sound enclosures
$\underline{f}\leq f\leq\overline{f}$ on $\calI_x$ for
$f\in\{F,\partial F/\partial x,\partial^2 F/\partial x^2\}$.
If~\eqref{eq:interval-bb} holds, then on~$\calI_x$:
$\partial F/\partial x\leq 0$, $\partial^2 F/\partial x^2\geq 0$, and $F>L_\infty$,
and $\lvert F\rvert<\infty$.
In particular, writing $F=G+g$ with $G=\Lhat^{(j-1)}$, conditions~(i)--(iv) and the finiteness half of~(vi) hold on the continuum configurations covered by the product of boxes~$\calI_x$ (a compact slice of~$\widetilde{\bbH}$).
When $L_\infty>0$, the floor test also yields~(iii).
Smoothness~(A5) on the box follows from the $C^\infty$ operator set once finiteness holds.
\end{lemma}

\begin{proof}
Soundness means every true value on~$\calI_x$ lies in the returned interval.
Hence $\overline{\partial F/\partial x}\leq 0$ forces $\partial F/\partial x\leq 0$ everywhere on~$\calI_x$;
$\underline{\partial^2 F/\partial x^2}\geq 0$ forces $\partial^2 F/\partial x^2\geq 0$;
and $\underline{F}>L_\infty$ forces $F>L_\infty$ on the box (hence also $F>0$ when $L_\infty>0$), which is~(iv) and yields~(iii).
Boundedness of~$\overline{F}$ is the finiteness claim of~(vi).
The argument is one-sided: failure of~\eqref{eq:interval-bb} does not refute~\eqref{eq:stage-admiss-bb}.
\end{proof}

\begin{lemma}[Structural exponent]\label{lem:ord}
Let $g$ be a finite expression tree over $\{+,-,\times,\div\}\cup\{\mathrm{pow}_p:p\in\calP\}$ in which the scale leaf~$x$ appears.
Define $\mathrm{ord}(g,x)$ recursively as in \cref{sec:stage-admiss}(b).
Then $g(\mathbf{h})|_{x=t}=\Theta\bigl(t^{\mathrm{ord}(g,x)}\bigr)$ as $t\to\infty$ (leading homogeneous term), and hence~\eqref{eq:stage-admiss-bb}(v) holds with $c_{x,j}=\mathrm{ord}(g,x)$ if and only if $\mathrm{ord}(g,x)<c_x^{(0)}$.
\end{lemma}

\begin{proof}
Proceed by induction on tree size.
Constants and leaves other than~$x$ are $O(1)=\Theta(t^{0})$.
A factor $\mathrm{pow}_p(x)$ contributes $\Theta(t^{p})$.
Products multiply leading powers (exponents add); quotients subtract; a sum is dominated by the child of largest exponent, so the inductive clause ``sums take the maximum'' records the leading order.
(If two children share the same maximal exponent their leading coefficients might cancel, in which case the true order is strictly smaller; the recursive $\mathrm{ord}$ is then a sound upper bound, and the test $\mathrm{ord}(g,x)<c_x^{(0)}$ remains sufficient for~(v).)
Thus $g|_{x=t}=O\bigl(t^{\mathrm{ord}(g,x)}\bigr)$, and whenever there is no exact leading-term cancellation one has matching $\Theta$ asymptotics.
Condition~(v) asks for some $c_{x,j}<c_x^{(0)}$ with $g=O(t^{c_{x,j}})$, which holds with $c_{x,j}=\mathrm{ord}(g,x)$ precisely when $\mathrm{ord}(g,x)<c_x^{(0)}$.
\end{proof}

\subsection{Zero-penalty implies certificates}\label{sec:soft-proof}
\begin{proof}[Proof of \cref{prop:soft}]
By definition each $v_a(g)\geq 0$, and $V(g)=\sum_a v_a(g)$.
Hence $V(g)=0$ if and only if every score vanishes.

\paragraph{Implication.}
$v_{\mathrm{mono}}(g)=0$ iff $\overline{\partial F/\partial x}\leq 0$ for all $x\in\{N,D\}$, which is the monotonicity half of~\eqref{eq:interval-bb}; likewise $v_{\mathrm{conv}}=0$ is equivalent to the convexity test and $v_{\mathrm{irred}}=0$ is equivalent to the floor test of \cref{lem:ia-sound} (and~(iii) when $L_\infty>0$).
The remaining interval test~$\overline{F}<\infty$ is not encoded as a soft score~$v_a$: under the smooth operator set of \cref{sec:boosting}, a completed DualInterval evaluation already yields finite enclosures, so~(A5) /~(vi) holds whenever the IA layer succeeds.
$v_{\mathrm{leaf}}(g)=0$ iff both scale variables appear as leaves.
$v_{\mathrm{decay}}(g)=0$ iff $\mathrm{ord}(g,x)\leq c_x^{(0)}-\varepsilon_{\mathrm{decay}}<c_x^{(0)}$, which is a strict strengthening of~(v).
Thus $V(g)=0$ together with finite DualInterval enclosures implies certificates~(a)--(b) (with the decay margin), and \cref{thm:guarantee} applies stage by stage.

\paragraph{Measured gaps.}
Soundness of DualInterval enclosures yields
$\partial F/\partial x\leq\overline{\partial F/\partial x}$ and
$\partial^2 F/\partial x^2\geq\underline{\partial^2 F/\partial x^2}$ on~$\calI_x$, so
\[
  \max\bigl(0,\partial F/\partial x\bigr)
  \;\leq\;
  \max\bigl(0,\overline{\partial F/\partial x}\bigr)
\]
pointwise on the box (and likewise for the convexity and floor gaps).
Taking maxima over $x\in\{N,D\}$ shows that $v_{\mathrm{mono}}$, $v_{\mathrm{conv}}$, and $v_{\mathrm{irred}}$ upper-bound the corresponding true violations on the compact slice of~$\widetilde{\bbH}$.
The quantities $\mathrm{ord}(g,x)$ and $\mathrm{leaves}(g)$ are exact combinatorial reads of the tree, so $v_{\mathrm{decay}}$ and $v_{\mathrm{leaf}}$ are exact.

\end{proof}

\subsection{Boosting guarantee}\label{sec:guarantee-proof}
\begin{proof}[Proof of \cref{thm:guarantee}]
Induct on the stage index~$j$.
For $j=0$, the claim is \cref{prop:chinchilla-admiss}: $\Lhat_0\in\calS$ with $L_\infty^{(0)}=E$ and the stated $c_x^{(0)}$.

Fix $j\geq 1$ and assume the claim holds through stage~$j-1$.
Certificates~(a)--(b) hold for $g=\Lhat_j$, so by \cref{lem:ia-sound} conditions~(i)--(iv) and~(vi) hold on each certification box~$\calI_x$ (hence~(A1)--(A5) on that compact slice of~$\widetilde{\bbH}$), and by \cref{lem:ord} condition~(v) holds on the half-line.
The box covers the observed grid; on $x>x_{\max}$,~(v) together with $G=\Lhat^{(j-1)}\in\calS$ implies that the correction is negligible relative to the leading excess of~$G$, so~(A1)--(A2) of~$G$ persist for large~$x$ (the derivatives of an $O(t^{c})$ term with $c<c_x^{(0)}<0$ are $o$ of those of the leading $t^{c_x^{(0)}}$ excess), and~$C^2$ regularity~(A5) extends by the operator set.
\Cref{prop:stage-iff} therefore yields $\Lhat^{(j)}\in\calS$ with the same joint floor $L_\infty^{(j)}=L_\infty^{(j-1)}=E$.
The marginal-floor / rate calculation in the proof of \cref{prop:stage-iff} gives~\eqref{eq:guarantee-bb}.
\end{proof}

\subsection{Stage-0 Chinchilla admissibility}\label{sec:chinchilla-proof}
\begin{proof}[Proof of \cref{prop:chinchilla-admiss}]
Write the hyperparameter remainder
\[
    R(\mathbf{h})
    \;\coloneqq\;
    \sum_{i=1}^{m} \frac{C_i}{\phi_i(h_i)^{\gamma_i}} + E,
\]
so $\Lhat_0 = A N^{-\alpha} + B D^{-\beta} + R(\mathbf{h})$.
Every summand of~$R$ is strictly positive on~$\widetilde{\bbH}$, hence $R(\mathbf{h}) > E > 0$.
Shape axioms~(A1),~(A2),~(A4), and~(A5) are imposed only along $x \in \{N,D\}$ with the other coordinates held fixed; under that restriction $R$ is constant in~$x$.

\paragraph{(A1)--(A2).}
Fix $x = N$ and freeze $(D,h_1,\ldots,h_m)$. Then
\[
    \frac{\partial \Lhat_0}{\partial N}
    \;=\;
    -\alpha A\, N^{-\alpha-1}
    \;<\; 0,
    \qquad
    \frac{\partial^2 \Lhat_0}{\partial N^2}
    \;=\;
    \alpha(\alpha+1)\, A\, N^{-\alpha-2}
    \;>\; 0
\]
for all $N \geq N_{\min} > 0$. The same identities with $(A,\alpha,N)$ replaced by $(B,\beta,D)$ give~(A1)--(A2) along~$D$.

\paragraph{(A3).}
Set $L_\infty \coloneqq E > 0$. For every continuum configuration with $N,D \geq$ their minima
(and hyperparameters in their~$H_h$),
\[
    \Lhat_0(\mathbf{h}) - E
    \;=\;
    \frac{A}{N^{\alpha}} + \frac{B}{D^{\beta}}
    + \sum_{i=1}^{m} \frac{C_i}{\phi_i(h_i)^{\gamma_i}}
    \;>\; 0,
\]
so the joint-floor inequality of~(A3) holds strictly on~$\widetilde{\bbH}$.
Along $N$ with the other coordinates fixed,
\[
    L^{\infty}_{N}
    \;=\;
    \lim_{t\to\infty} \Lhat_0\big|_{N=t}
    \;=\;
    \frac{B}{D^{\beta}} + R(\mathbf{h})
    \;=\;
    \frac{B}{D^{\beta}}
    + \sum_{i=1}^{m} \frac{C_i}{\phi_i(h_i)^{\gamma_i}} + E,
\]
which is finite and satisfies $L^{\infty}_{N} \geq E = L_\infty$ (and is strictly larger whenever $B>0$ or $m\geq 1$).
The marginal floor $L^{\infty}_{D}$ is analogous.

\paragraph{(A4).}
With the same frozen coordinates,
$\Lhat_0\big|_{N=t} - L^{\infty}_{N} = A\, t^{-\alpha}$, hence
\[
    \lim_{t\to\infty}
    \frac{\log\bigl(\Lhat_0\big|_{N=t} - L^{\infty}_{N}\bigr)}{\log t}
    \;=\; -\alpha \;=\; c_N^{(0)} \;<\; 0.
\]
The $D$-axis calculation yields $c_D^{(0)} = -\beta < 0$.

\paragraph{(A5).}
On $[N_{\min},\infty)$ the map $N \mapsto A N^{-\alpha}$ is $C^\infty$ (hence $C^2$), the remaining terms are independent of~$N$, and
$\Lhat_0(N_{\min},\cdot) = A N_{\min}^{-\alpha} + B D^{-\beta} + R(\mathbf{h})$ is finite
(and strictly larger than~$E$ by~(A3)).
The same holds with $N$ replaced by~$D$.

Thus $\Lhat_0$ satisfies~(A1)--(A5), i.e.\ $\Lhat_0 \in \calS$.
\end{proof}

\begin{remark*}
If the hyperparameter block in~\eqref{eq:chinchilla-bb} is replaced by nonnegative quadratic penalties in
$\log\phi_{\mathrm{lr}}(\mathrm{lr})$ and $\log\phi_{\mathrm{wd}}(\mathrm{wd})$
(as in the 4D U-shaped baseline of \cref{sec:boosting}), those terms remain constant along $N$ and~$D$.
Whenever $A,B,E,\alpha,\beta > 0$, the same derivative and asymptotic calculations apply verbatim, so~(A1)--(A5) continue to hold with the same $L_\infty = E$ and $c_x^{(0)}$.
\end{remark*}

\end{document}
